\documentclass[11pt]{article}
\usepackage[margin=25mm]{geometry}
\usepackage{amsmath,amssymb,booktabs}
\usepackage{eqannotate}
\title{EqAnnotate Article-Flow Integration Fixture}
\author{Integration Test}
\date{}
\newcommand{\demotext}{Modern scientific papers often mix prose, displayed mathematics, cross-references, tables, and figures on the same page. EqAnnotate keeps annotation space attached to the equation while preserving document flow, readable nearby paragraphs, ordinary numbering, and references.}
\begin{document}
\maketitle
\begin{abstract}
This synthetic article checks EqAnnotate inside continuous prose. It combines automatic and manual annotations, numbered displays, long labels, multiple displays, floats, and multi-line mathematics.
\end{abstract}

\section{Introduction}
\demotext

\demotext

The first annotated display is referenced as Eq.~\eqref{eq:flow}. Text immediately before the display should remain visually separated from the upper callouts.
\begin{annotatedequation}[numbered]
\frac{dx_t}{dt}=\eqmark[blue]{velocity}{v_\theta(x_t,t)}+\eqmark[orange]{guidance}{\lambda g(x_t,t)}
\label{eq:flow}
\eqannotate{velocity}{Learned velocity field}
\eqannotate{guidance}{Guidance correction}
\end{annotatedequation}
The package-managed lower band separates the following paragraph from the callouts and keeps page spacing cohesive. \demotext

\begin{figure}[t]
\centering
\rule{0.82\linewidth}{24mm}
\caption{A placeholder float used to check interaction between normal page layout and annotated displays.}
\label{fig:placeholder}
\end{figure}

\section{Mixed Automatic and Manual Placement}
Figure~\ref{fig:placeholder} is a normal float. The next equation keeps two terms under automatic control and places one callout manually.
\begin{annotatedequation}[numbered]
\mathcal{L}=\eqmark[blue]{data}{\mathcal{L}_{\mathrm{data}}}
+\eqmark[orange]{reg}{\lambda\mathcal{L}_{\mathrm{reg}}}
+\eqmark[green]{aux}{\gamma\mathcal{L}_{\mathrm{aux}}}
\label{eq:loss}
\eqannotate{data}{Data fidelity}
\eqannotate{reg}{Regularizer}
\eqannotatemanual[xshift=25mm,yshift=-16mm]{aux}{Manual auxiliary term}
\end{annotatedequation}
The manual label participates in document-space reservation, so surrounding article text follows it without an additional spacing command.

\begin{table}[b]
\centering
\begin{tabular}{lcc}
\toprule
Setting & Automatic & Manual \\
\midrule
Placement & yes & explicit \\
Space reservation & yes & yes \\
Theme inheritance & yes & yes \\
\bottomrule
\end{tabular}
\caption{Expected behavior of the two placement paths.}
\end{table}

\demotext

\section{Consecutive Displays and Long Labels}
Two annotated displays may occur with only a short sentence between them. Their reserved bands remain distinct with controlled vertical spacing.
\begin{annotatedequation}
p(x)=\frac{1}{\sqrt{2\pi}\eqmark[blue]{sigma}{\sigma}}
\exp\left(-\frac{(x-\eqmark[yellow]{mu}{\mu})^2}{2\sigma^2}\right)
\eqannotate{mu}{Mean}
\eqannotate{sigma}{Scale}
\end{annotatedequation}
A short bridge sentence separates the displays.
\begin{annotatedequation}
\hat y=\eqmark[blue]{predictor}{f_\theta(x)}+\eqmark[orange]{residual}{r_\phi(x)}
\eqannotate{predictor}{Primary prediction produced by the learned model before the residual correction is applied}
\eqannotate{residual}{Residual correction}
\end{annotatedequation}
The long annotation wraps within the active text width, and the following prose begins below its measured label. \demotext

\section{Multi-line Mathematics}
The aligned wrapper is tested inside ordinary prose.
\begin{annotatedalign}[numbered]
r_0 &= \eqmark[blue]{initial}{b-Ax_0},\\
p_k &= M^{-1}r_k,\\
x_{k+1} &= x_k+\eqmark[orange]{step}{\alpha_k p_k}
\label{eq:iter}
\eqannotate{initial}{Initial residual}
\eqannotate{step}{State update}
\end{annotatedalign}
Equation~\eqref{eq:iter} has one number for the complete aligned block. \demotext

\section{Conclusion}
\demotext

\end{document}
